\documentclass[12pt,a4paper,oneside]{article}

\usepackage[utf8]{inputenc}
\usepackage[T1]{fontenc}
\usepackage[brazil]{babel}
\usepackage{geometry}
\geometry{a4paper, margin=2.5cm}
\usepackage{indentfirst}
\usepackage{times}
\usepackage{graphicx}
\usepackage{hyperref}
\usepackage{abntex2cite}

\setlength{\parindent}{1.25cm}
\setlength{\parskip}{0.3cm}

\begin{document}

\begin{center}
  \textbf{\large ESCOLA E FACULDADE DE TECNOLOGIA SENAI GASPAR RICARDO JÚNIOR}\\
  \textbf{Desenvolvimento de Sistemas}\\[1cm]
  \textbf{\Large Lucas Miguel Leite}\\[0.5cm]
  \textbf{Professor: Deivison Takatu}\\[2cm]
  \textbf{\LARGE\textbf{Impactos da TI nas Organizações: Análise Transversal das Perspectivas da PwC e McKinsey}}\\[1cm]
  Sorocaba -- 11/03/2026
\end{center}

\vspace{1cm}

\begin{center}
  \textbf{RESUMO}
\end{center}

Este artigo analisa os impactos da Tecnologia da Informação nas organizações a partir de uma perspectiva transversal dos relatórios da PwC e da McKinsey, evidenciando três dimensões fundamentais interconectadas: Dados como fundação tecnológica, Estratégia como direcionamento para o valor (ROI), e Impactos Organizacionais na força de trabalho e na cultura. A PwC enfatiza a Inteligência Artificial Generativa como ferramenta de modernização de dados, a abordagem de portfólio para adoção em escala e a emergência da figura do ``generalista em IA''. A McKinsey foca na infraestrutura de conectividade (5G/IoT, Cloud, Edge Computing), na convergência de tendências tecnológicas e na criação de culturas ágeis para o futuro do trabalho. Conclui-se que a falha em conectar qualquer um dos três pontos (Estratégia, Dados, Organização) resulta em investimentos sem retorno e perda de vantagem competitiva.

\vspace{0.5cm}
\textbf{Palavras-chave:} Impactos da TI. Inteligência Artificial. Transformação Digital. PwC. McKinsey.

\newpage

\section{Introdução}

Tanto a PwC quanto a McKinsey demonstram que a tecnologia deixou de ser uma área de suporte operacional para se tornar o núcleo da vantagem competitiva. Para descrever os impactos da TI, os relatórios conectam três dimensões fundamentais: Dados, Estratégia e Organização, formando um modelo integrado no qual a falha em qualquer um dos pontos compromete toda a cadeia de valor \cite{tarapanoff2015}.

\section{Dados: A Fundação Tecnológica}

Os dados são tratados como a matéria-prima essencial para qualquer transformação. O impacto da TI começa na capacidade de capturar, organizar e proteger essas informações.

\subsection{PwC: IA e Dados Não Estruturados}

A PwC destaca que a Inteligência Artificial Generativa (GenAI) é a principal ferramenta para modernizar a estratégia de dados das empresas, permitindo extrair valor de dados não estruturados (documentos, logs, apresentações). O diferencial competitivo não está em qual modelo de IA a empresa usa, mas em como ela aplica a IA sobre seus dados proprietários \cite{pwc2024}.

\subsection{McKinsey: Infraestrutura e Conectividade}

A McKinsey foca na arquitetura necessária para movimentar dados. Tendências como o Futuro da Conectividade (5G/IoT) e a Infraestrutura Distribuída (Cloud e Edge Computing) permitem processamento em tempo real. Destacam-se ainda as ``Arquiteturas de Confiança'' (Zero Trust e Blockchain) para garantir segurança, privacidade e governança dos dados em um mundo hiperconectado.

\section{Estratégia: Direcionamento para o Valor}

A adoção tecnológica baseada apenas no ``hype'' falha. Ambas as consultorias enfatizam que a tecnologia deve ser orientada por estratégias de negócios top-down.

\subsection{PwC: Portfólio e IA Responsável}

A PwC sugere migração de projetos piloto isolados para adoção em escala, com abordagem de portfólio que equilibra ganhos rápidos (ground game) com inovações disruptivas (moonshots). A estratégia de TI está ligada à IA Responsável (governança, ética e gestão de riscos), vista como requisito fundamental para garantir ROI e confiança dos stakeholders.

\subsection{McKinsey: Convergência e Aplicação Prática}

A estratégia vencedora combina múltiplas tendências (IA Aplicada + Automação de Processos + Cloud). A McKinsey enfatiza a ``Industrialização do Machine Learning'' e o alinhamento das iniciativas tecnológicas diretamente à cadeia de valor, otimizando eficiência operacional e criando novos modelos de receita \cite{rocha2018}.

\section{Impactos Organizacionais: Pessoas e Cultura}

O impacto da TI culmina na transformação da força de trabalho, exigindo novos modelos mentais e estruturas corporativas.

\subsection{PwC: A Era dos Agentes de IA}

A automação de tarefas rotineiras e a adoção de Agentes de IA mudarão a estrutura da força de trabalho, assumindo formato de ampulheta (foco em juniores apoiados por IA e seniores focados em estratégia). Surgirá a figura do ``generalista em IA'' --- profissionais que orquestram e auditam o trabalho das máquinas, exigindo forte colaboração multidisciplinar \cite{pwc2024}.

\subsection{McKinsey: Cultura Ágil e Futuro do Trabalho}

A tecnologia impacta diretamente a produtividade em funções baseadas no conhecimento. Com inovações como o Software 2.0 (IA escrevendo código), o foco humano desloca-se da execução técnica para resolução de problemas complexos. A adoção exige criação de cultura ágil e estratégias de retenção de talentos especializados.

\section{Conclusão}

A análise transversal da PwC e McKinsey revela que os impactos da TI em uma corporação moderna operam em três vértices interdependentes. A Estratégia define o problema de negócio e garante o uso ético das ferramentas; os Dados, suportados por infraestrutura de ponta e conectividade, alimentam os algoritmos e sistemas; e a Organização, com cultura ágil e novos papéis de trabalho, executa a transformação lado a lado com sistemas inteligentes. A falha em conectar qualquer um desses três pontos resulta em investimentos sem retorno e perda de vantagem competitiva \cite{kluyver2007}.

\bibliographystyle{plain}
\bibliography{referencias}

\end{document}
